\textbf{Predstavil:}

V implicitni ali polarni obliki so podane naslednje krivulje:
		\begin{itemize}
				% \item $ r(\varphi) = \frac{a}{\varphi}, \quad \varphi \in \mathbb{R}^+$, (hiperbolična spirala).
				% \item $r(\varphi) = e^{a\varphi}$, (logaritemska spirala).
				\item $r(\varphi) = \cos(3\varphi), \quad \varphi \in \left[-\frac{\pi}{6}, \frac{\pi}{6} \right] \cup \left[\frac{\pi}{2}, \frac{5\pi}{6} \right] \cup \left[\frac{7\pi}{6}, \frac{9\pi}{6} \right]$.
				\item $r(\varphi) = a(1 + \cos(\varphi)), \quad \varphi \in [-\pi, \pi] $, (kardioida).
				% \item $(x^2 + y^2) = 2a^2(x^2 - y^2), \quad a > 0$, (lemniskata). \textvf{Nasvet}: Vpelji polarne koordinate.
		\end{itemize}
		\begin{enumerate}
		\item Čimbolj natančno narišite dane krivulje. Nasvet: skicirajte grafe $r(\varphi)$, $x(\varphi)$ in $y(\varphi)$; poskusite določiti ekstremne točke, t.j. rešite enačbi $\dot{x}=0$ in $\dot{y}=0$.
		\item Izračunajte naklon krivulje oziroma $y' =\frac{\dot{y}}{\dot{x}}$. Določite tangento na dano krivulje v točki, ko je kot $\varphi = \frac{\pi}{2}$, če obstaja.
		\end{enumerate}
